\bookchapter{פרק 4: התקנה וניהול של אפליקציות — “ערכת הכישורים” של הטלפון}

\chapterintro{אפליקציות הופכות את הטלפון לכל-יכול, אבל צריך גם לנהל אותן נכון. בדיוק כפי שמסדרים לעובד חדש שולח ומקצים הרשאות, גם התקנת אפליקציות דורשת תהליך מסודר.}

\booksection{החנות הרשמית: “סופר האפליקציות” הבטוח ביותר}

\begin{itemize}
  \item \textbf{אנדרואיד}: Google Play (בחו"ל) או החנות המובנית של המותג (כגון Huawei AppGallery, Xiaomi App Store).
  \item \textbf{iOS}: רק App Store.
  \item \textbf{HarmonyOS}: Huawei AppGallery.
\end{itemize}

המשמעות הביטחונית המרכזית: אפליקציות שעברו אישור, עדכונים מובטחים ותשלומים בטוחים. קו מנחה:\textbf{הורידו אפליקציות רק מחנויות רשמיות!}

\booksection{מצב טהור/מצב הגנה: הפעלת “חומת ההגנה” של הטלפון}

מותגים סיניים מובילים (Huawei, Xiaomi, OPPO, vivo, Honor) כוללים מובנה\textbf{מצב טהור}או\textbf{מצב הגנה}— אלה מתגי אבטחה ברמת המערכת.\textbf{מומלץ בחום לכל המשתמשים להפעיל אותם}.

\booksubsection{מה זה בעצם?}

ה\textbf{פונקציה המרכזית}של מצב טהור (הנקרא גם “מצב הגנה”) היא\textbf{חסימת התקנת אפליקציות צד שלישי}: כשהוא מופעל, המערכת\textbf{מאפשרת התקנת אפליקציות רק מהחנות הרשמית}; חבילות התקנה שהורדו מדפדפן, מקוד QR, מתוכנות צ'אט וכדומה נחסמות ישירות — מניעת תוכנות זדוניות עוד במקור.

\booksubsection{למי זה מתאים?}

\begin{itemize}
  \item \textbf{משתמשים מבוגרים}: לא מבחינים בין פרסומות "רשמיות" ל"מזויפות" — מצב טהור הוא קמע ההגנה הבטוח ביותר.
  \item \textbf{משתמשים שאינם רגישים להרשאות אפליקציות}: לא צריכים להתקין אפליקציות נישה, גרסאות בטא או משחקים משרתים זרים.
  \item \textbf{כל המשתמשים שרוצים לשפר את אבטחת המכשיר}: אפילו לחובבי טכנולוגיה — עדיף מופעל מאשר כבוי; אם יש צורך זמני, פשוט כבו אותו לרגע.
\end{itemize}

\booksubsection{נתיב ההפעלה (בדוגמת Huawei/Xiaomi)}

\begin{itemize}
  \item \textbf{Huawei/HarmonyOS}: [הגדרות] > [מערכת ועדכון] >\textbf{[מצב טהור]}, בחרו "הפעל".
  \item \textbf{Xiaomi}: [הגדרות] > [הגנת פרטיות] >\textbf{[הגן על הפרטיות]}, או [הגדרות] > [הגדרות אפליקציות] > [התקנת אפליקציות] — כבו את "התקן אפליקציות לא מוכרות".
  \item \textbf{OPPO/vivo}: [הגדרות] > [אבטחה] >\textbf{[התקנת אפליקציות ממקורות חיצוניים]}— כבו את ההרשאות אחת-אחת.
\end{itemize}

\booksubsection{מנגנון בדיקת האבטחה של חנות האפליקציות}

חנות האפליקציות אינה “קישוט” — מאחוריה פועל צינור אבטחה שלם:

\begin{itemize}
  \item \textbf{סריקת וירוסים}: לפני העלאתה, כל אפליקציה עוברת סריקה משותפת של מספר מנועי אנטי-וירוס, המזהה סוסים טרויאניים ווירוסים מוכרים.
  \item \textbf{בדיקת הרשאות}: החנות בודקת אם ההרשאות שהאפליקציה מבקשת "הגיוניות" — אפליקציית פנס שדורשת גישה לאנשי קשר? נדחית או מורדת מהמדף מיד.
  \item \textbf{אימות מפתחים}: העלאה מחייבת רישיון עסק, תעודת זהות אישית או אימות כרטיס אשראי. אם משהו משתבש — אפשר למצוא את האחראי; זו לא "אפליקציה צללית".
  \item \textbf{ניטור רציף}: גם אפליקציות שכבר עלו לאוויר נסרקות מחדש באופן רציף; התנהגות זדונית שמתגלה מביאה להורדה כפויה מהמדף.
\end{itemize}

\booksection{מקורות צד שלישי ו“טפילים אלקטרוניים”: סיכון גבוה ביותר!}

הורדת אפליקציות מדפי אינטרנט או מקישורים לא מוכרים עלולה להכניס וירוסים, לגנוב פרטיות ולבזוק כספים מהתשלומים. כאן נציג מושג חדש:\textbf{“טפילים אלקטרוניים”}.

\textbf{מהם “טפילים אלקטרוניים”?}\par

הכוונה ל\textbf{תוכנות “גרסה מפוצחת” שהושתל בהן סוס טרויאני}— אתה חושב ש“חסכת כסף”, ובפועל “נפלת בפח”. צורות נפוצות:

\begin{itemize}
  \item תוכנות בתשלום ב"גרסה מפוצחת", "גרסה חינמית" או "גרסה ללא פרסומות";
  \item כלים אפורים כמו "הגדלת לייקים", "חטיפת מעטפות אדומות" ו"צ'יטים למשחקים";
  \item אפליקציות "גרסה פנימית" כביכול, המומלצות ב"פרסומות קטנות" בצ'אטים קבוצתיים.
\end{itemize}

התוכנות האלה נראות תקינות על פני השטח, אך מאחורי הקלעים עלולות\textbf{לגנוב אנשי קשר, לקרוא SMS, לצותת למיקרופון ולכרות מטבעות ברקע}. מצב טהור הוא\textbf{קו ההגנה הראשון}מפניהן — הוא פשוט לא מאפשר לך להתקין אותן.

\booksubsection{הצעות לקו הגנה השני:}

\begin{itemize}
  \item \textbf{סריקת וירוסים מלאה תקופתית}: השתמשו במרכז האבטחה המובנה של הטלפון (כגון Huawei Phone Manager, Xiaomi Security Center) לסריקה פעם בחודש.
  \item \textbf{הענקת הרשאות רגישות בזהירות}: אנשי קשר, SMS, מיקום, מצלמה, מיקרופון — חמש "ההרשאות המסוכנות" האלה חייבות אישור קפדני.
  \item \textbf{שימו לב לצריכת חשמל ולתעבורת נתונים חריגות}: הטלפון מתחמם פתאום, הסוללה מתרוקנת מהר, תעבורת נתונים בלילה — כנראה מסתתר "טפיל" ברקע.
\end{itemize}

אם בכל זאת אין ברירה וצריך להתקין אפליקציית צד שלישי:

\begin{enumerate}
  \item ודאו שהמקור אמין לחלוטין.
  \item בדקו היטב את דרישות ההרשאות לפני ההתקנה.
  \item סרקו עם אפליקציית אבטחה לאחר מכן.
\end{enumerate}

ההמלצה הסופית: בהגדרות המערכת, כבו את ההרשאה “התקן אפליקציות לא מוכרות”.

\booksection{התקנה, עדכון והסרה}

\begin{itemize}
  \item \textbf{התקנה}: לחצו "קבל" או "התקן" בחנות; בפתיחה הראשונה לאחר ההתקנה יבקשו ממך הרשאות.
  \item \textbf{עדכון}:\textbf{חובה לעדכן!}בעיקר לתיקון פרצות אבטחה. מומלץ להפעיל עדכונים אוטומטיים.
  \item \textbf{הסרה}: לחיצה ארוכה על האייקון וגרירה לאזור ההסרה. רוצה לנקות בצורה יסודית יותר? לך ל-[הגדרות] > [ניהול אפליקציות], מצא והסר.
\end{itemize}

\booksection{ניהול הרשאות: אתה “הבוס” של האפליקציות}

אפליקציות דורשות הרשאות (כמו אנשי קשר, מיקום) — לפעמים זה הגיוני (מפות צריכות מיקום), לפעמים לא הגיוני (פנס דורש אנשי קשר).

\begin{itemize}
  \item \textbf{אפשרויות הרשאה}: "אפשר תמיד", "אפשר רק בזמן שימוש", "חסום", "שאל בכל פעם".
  \item \textbf{עקרונות ניהול}:
  \begin{enumerate}
    \item מינימיזציה: אל תיתן אם זה לא הכרחי.
    \item ביקורת לאחר מעשה: פעם בתקופה עבור על הגדרות ההרשאות ובדוק מי "חורג מסמכותו".
    \item הקשר: כמו אפליקציית מוניות, תן הרשאת מיקום רק בזמן השימוש.
  \end{enumerate}
\end{itemize}

\begin{tcbwarn}{הרשאות שדורשות פיקוח צמוד}
\begin{itemize}
  \item אנשי קשר/SMS: עלולים לשמש לשליחת הודעות זבל.
  \item מיקום בזמן אמת: עלול לשמש למעקב אחר מיקומך.
  \item מצלמה/מיקרופון: עלולים לשמש לצילום והקלטה בסתר.
  \item גישה לקבצים: עלולה לקרוא את הקבצים הפרטיים שלך.
\end{itemize}

\end{tcbwarn}

\booksection{מבנה קובץ APK של Android}

APK(Android Package Kit)הוא בעצם קובץ ZIP דחוס,לאחר פריסה ניתן לראות את התיקיות והקבצים המרכזיים הבאים:

\booksubsection{AndroidManifest.xml}

:“תעודת הזהות” ו“ספר ההוראות” של האפליקציה.מתעד שם חבילה,מספר גרסה,הרשאות נדרשות,גרסת מערכת נתמכת,Activity כניסה ומידע מרכזי נוסף.המערכת קוראת קובץ זה ראשון בעת התקנת האפליקציה.

\begin{itemize}
  \item \textbf{classes.dex}:קוד הליבה של האפליקציה.קוד מקור Java/Kotlin מהודר ל-Dalvik bytecode,מאוחסן בקובץ זה.לאפליקציה עשויים להיות מספר קבצי dex(כגון classes2.dex, classes3.dex).
  \item \textbf{resources.arsc}:טבלת אינדקס משאבים מהודרת.מתעדת את כל מזהי המשאבים והקשרים ביניהם.
  \item \textbf{תיקיית res/}:תיקיית משאבים.מכילה תמונות,קבצי פריסה,מחרוזות,הגדרות צבע ועוד.תמונות ברזולוציות שונות ממוקמות בתיקיות drawable שונות(כגון drawable-hdpi, drawable-xhdpi).
  \item \textbf{תיקיית assets/}:תיקיית משאבים גולמיים.מאחסנת משאבים שאינם עוברים הידור,כגון קבצי דפי אינטרנט,גופנים,שמע ועוד.ניתן לקרוא בזמן ריצה דרך AssetManager.
  \item \textbf{תיקיית lib/}:תיקיית ספריות מקומיות.מאחסנת קבצי .so מהודרים(תוצרי הידור קוד C/C++),מסווגים לפי ארכיטקטורת CPU(כגון arm64-v8a, armeabi-v7a).משחקים ואפליקציות הדורשות חישוב בביצועים גבוהים בדרך כלל כוללות תיקייה זו.
  \item \textbf{תיקיית META-INF/}:תיקיית מידע חתימה.מכילה שלושה קבצים MANIFEST.MF, CERT.SF, CERT.RSA,לצורך אימות שלמות וחתימת APK.
  \item \textbf{2. מבנה קובץ IPA של iOS}IPA(iOS App Store Package)הוא גם קובץ דחוס,המבנה שונה מעט:
\end{itemize}

\booksubsection{תיקיית Payload/}

:מיקום התוכן המרכזי.בתוכה יש תיקיית .app,המכילה את כל תוכן האפליקציה.

\begin{itemize}
  \item \textbf{Info.plist}:קובץ התצורה של האפליקציה,דומה ל-AndroidManifest.xml של Android.מתעד מזהה אפליקציה,גרסה,דרישות הרשאות,מכשירים נתמכים ועוד.
  \item \textbf{קובץ בר-הרצה}:התוכנית הראשית בעלת שם זהה לתיקיית .app,היא קובץ בינארי מהודר בפורמט Mach-O.
  \item \textbf{תיקיית Frameworks/}:ספריות דינמיות ומסגרות תלויות של האפליקציה.
  \item \textbf{קבצי משאבים}:תמונות,שמע,וידאו ומשאבים אחרים,מאוחסנים בדרך כלל בתוך תיקיית .app.
  \item \textbf{תיקיית \_CodeSignature/}:מידע חתימת קוד,לצורך אימות שלמות ומקור האפליקציה.
  \item \textbf{3. מנגנון אימות חתימה}למה אפליקציות מחנויות רשמיות בטוחות יותר?אימות חתימה הוא המפתח.
\end{itemize}

\booksubsection{חתימת Android}

:המפתח חותם על APK במפתח פרטי,המערכת מאמתת במפתח ציבורי בעת ההתקנה.אם APK שונה(למשל הוחדר וירוס),החתימה תיפסל,המערכת תסרב להתקין.Google Play גם דורש App Signing, Google שומרת על מפתחות החתימה,מגנה עוד יותר על אבטחת האפליקציה.

\begin{itemize}
  \item \textbf{חתימת iOS}:הביקורת של אפל מחמירה יותר.כל IPA חייב להיות חתום בתעודה שהונפקה על ידי אפל,אחרת לא ניתן להתקין.זו הסיבה ש-iOS יכולה להתקין אפליקציות רק מ-App Store או מתעודות ארגוניות.לאפליקציות חתומות בתעודה ארגונית יש תוקף מוגבל,ואפל יכולה לבטל בכל עת.
  \item \textbf{תפקידי החתימה}:
  \begin{enumerate}
    \item אימות מקור: וידוא שהאפליקציה אכן מגיעה מהמפתח המוצהר.
    \item אימות שלמות: וידוא שהאפליקציה לא הודבקה במהלך ההעברה.
    \item מניעת התחזות: אפליקציה זדונית לא יכולה להתחזות לאפליקציה אחרת.
  \end{enumerate}
\end{itemize}

\booksubsection{גודל חבילת ההתקנה ואופטימיזציה}

ייתכן שתגלה שהגודל המוצג של אפליקציה בחנות שונה מהשטח שהיא תופסת לאחר ההתקנה. הסיבה לכך:

\begin{itemize}
  \item \textbf{גודל ההורדה}: גודל ה-APK/IPA בדחיסה.
  \item \textbf{גודל ההתקנה}: הגודל לאחר פריסה, בדרך כלל גדול בהרבה מגודל ההורדה.
  \item \textbf{תפוסה בזמן ריצה}: בזמן ריצה האפליקציה מייצרת גם מטמון וקבצי נתונים, ותופסת עוד יותר מקום.
\end{itemize}

שיטות האופטימיזציה של אפליקציות מודרניות:

\begin{itemize}
  \item \textbf{App Bundle (אנדרואיד)}: טכנולוגיית הפצה דינמית של Google Play. בזמן ההורדה המשתמש מקבל רק את המשאבים המתאימים למכשירו (למשל הורדת תמונות רק ברזולוציית המסך המתאימה), מה שמפחית משמעותית את גודל ההורדה.
  \item \textbf{App Thinning (iOS)}: טכנולוגיה דומה; אפל תקצץ אוטומטית משאבים מיותרים בהתאם למכשיר.
  \item \textbf{הורדה לפי דרישה}: אפליקציות גדולות כמו משחקים יכולות להוריד תחילה את חבילת הליבה, ואת חבילות המשאבים מאוחר יותר.
\end{itemize}

\begin{tcbinfo}{טיפ: למה חבילות ההתקנה של חלק מהאפליקציות גדולות במיוחד?}
אפליקציות משחק בדרך כלל הגדולות ביותר, כיוון שהן מכילות טקסטורות HD רבות, מודלים תלת-ממדיים וקבצי שמע. אפליקציות חברתיות גם גדולות יחסית, כיוון שהן צריכות לתמוך במודולים פונקציונליים רבים. אפליקציות כלי עבודה בדרך כלל הקטנות ביותר. אם אתה מוצא אפליקציית כלי עבודה גדולה במיוחד, היזהר שמא היא מצרפת פונקציות מיותרות.

\end{tcbinfo}

\begin{tcbpractice}{תרגול מעשי: בדוק את ההרשאות של 3 אפליקציות}
בחר 3 אפליקציות נפוצות (כגון חברתית, קניות, כלי עבודה), היכנס לדף ניהול הרשאות האפליקציות ב-[הגדרות], בדוק פריט אחר פריט, כבה הרשאות רגישות מיותרות או שנה ל“אפשר רק בזמן שימוש”.

\end{tcbpractice}
